% TUFTE-BOOK CLASS {{{

%%%%%%%%%%%%%%%%%%%%%%%%%%%%%%%%%%%%%%%%%
% Tufte-Style Book (Documentation Template)
% LaTeX Template
% Version 1.0 (5/1/13)
%
% This template has been downloaded from:
% http://www.LaTeXTemplates.com
%
% Original author:
% The Tufte-LaTeX Developers (tufte-latex.googlecode.com)
%
% Edited by:
% Federico Marotta (federicomarotta@mail.com)
%
% License:
% Apache License (Version 2.0)
%
% IMPORTANT NOTE:
% In addition to running BibTeX to compile the reference list from the .bib
% file, you will need to run MakeIndex to compile the index at the end of the
% document.
%
%%%%%%%%%%%%%%%%%%%%%%%%%%%%%%%%%%%%%%%%%

% This document uses the tufte-book class (which in turn uses the 
% tufte-common class)
%
% Summary of available options for the tufte-book class:
% a4paper: use the A4 paper size
% sfsitednotes: use a sans serif font for sidenotes and title blocks
% twoside: print the page number on the outside edge
% symmetric: typeset the sidenotes on the outside edge instead of on the right side
% justified: justify text
% notitlepage: have \maketitle generate a title block instead of a title page
% notoc: suppress tufte's custom table of contents
% nohyper: do not load the hyperref package
% nobib: do not load natbib (which is incompatible with biblatex)
\documentclass[a4paper,oneside,marginals=justified,nobib]{tufte-book} 

% }}}

% PACKAGES AND OTHER DOCUMENT CONFIGURATIONS {{{
% --------------------------------------------------------------------------

% Encoding and fonts
\usepackage[utf8]{inputenc} % Use utf8 encoding in the input (.tex) file
\usepackage[english]{babel} % Load characters and hyphenation
\usepackage[T1]{fontenc} % Use utf8 encoding in the output (.pdf) file
\usepackage{hyphenat} % Enable hyphenation for special fonts
\usepackage[english=british]{csquotes} % Quote like a boss

% Special contents
\usepackage{textgreek} % Greek characters
\usepackage{amsmath} % Improved mathematics
\usepackage{amsfonts} % Math fonts
\usepackage[scale=2]{ccicons} % Creative Commons icons
\usepackage{units} % Used for printing standard units
\usepackage{xspace} % Used for printing a trailing space better than using a tilde (~) using the \xspace command
\usepackage{todonotes} % Add useful notes in the margins
\usepackage{subfiles} % Adopt a modular structure
\usepackage[most]{tcolorbox} % Make fancy boxes

% TOC
\setcounter{secnumdepth}{-1} % Do not number headings
\setcounter{tocdepth}{1} % Show headings up to sections in the toc
\usepackage[notlot,notlof]{tocbibind}
%\makeatletter
%\renewcommand*\l@chapter{\@dottedtocline{1}{0em}{2.3em}}
%\renewcommand*\l@section{\@dottedtocline{1}{0em}{2.3em}}
%\makeatother

% Images
\usepackage{graphicx} % Needed to insert images into the document
\setkeys{Gin}{width=\linewidth,totalheight=\textheight,keepaspectratio} % Improves figure scaling
\graphicspath{{images/}{../images/}} % Sets the default location of pictures
% (the second path is relative to the subfiles)

% Bibliography
% See https://github.com/Tufte-LaTeX/tufte-latex/issues/60 for the 
% problems concerning biblatex
\usepackage[citestyle=authoryear-icomp,bibstyle=authoryear,sorting=none,backend=biber]{biblatex}
\addbibresource{main.bib}
\AtEveryBibitem{\clearfield{url}} % Do not show URL
\renewcommand{\cite}[2][0pt]{\sidenote[][#1]{\fullcite{#2}} }

% Index
%\usepackage{makeidx} % Used to generate the index
%\makeindex % Generate the index which is printed at the end of the 
%document (remember to uncomment \printindex if you use the index)

% Document appearance
\usepackage{microtype} % Improves character and word spacing
\usepackage{booktabs} % Better horizontal rules in tables
\usepackage{fancyvrb} % Allows customization of verbatim environments
\usepackage{multicol}
\hypersetup{colorlinks} % Comment this if you don't want colored links
\fvset{fontsize=\normalsize} % Change here the font size of all verbatim 
\setlength{\headheight}{24pt} % Set the height of headers

% Print the chapter number and title on the same line 
% https://latex.org/forum/viewtopic.php?t=25750
\titleformat{\chapter}%
[block]% shape
{\relax\ifthenelse{\NOT\boolean{@tufte@symmetric}}{\begin{fullwidth}}{}}% format applied to label+text
{\itshape\huge\thechapter\hspace{10pt}}% label
{0pt}% horizontal separation between label and title body
{\huge\rmfamily\itshape}% before the title body
[\ifthenelse{\NOT\boolean{@tufte@symmetric}}{\end{fullwidth}}{}]% after the title body

% CUSTOM COMMANDS AND ENVIRONMENTS {{{
% --------------------------------------------------------------------------

% Document-specific {{{

% Print latin words in italics
\newcommand{\invitro}{\textit{in vitro}\xspace}
\newcommand{\invivo}{\textit{in vivo}\xspace}
\newcommand{\cis}{\textit{cis}\xspace}
\newcommand{\trans}{\textit{trans}\xspace}
\newcommand{\etal}{\textit{et al.}\xspace}
\newcommand{\etcetera}{\textit{et cetera}\xspace}
\newcommand{\etc}{\textit{etc.}\xspace}

% Easily label and reference elements
\newcommand{\labch}[1]{\label{ch:{#1} }}
\newcommand{\refch}[1]{Chapter \ref{ch:{#1} }}
\newcommand{\labsec}[1]{\label{sec:{#1} }}
\newcommand{\refsec}[1]{Section \ref{sec:{#1} }}
\newcommand{\labfig}[1]{\label{fig:{#1} }}
\newcommand{\reffig}[1]{Figure \ref{fig:{#1} }}

% Environment to put external abstracts in a box
\newenvironment{external_abstract}[1]
{
	\begin{fullwidth}
	\begin{tcolorbox}[#1]
	\footnotesize
}
{
	\end{tcolorbox}
	\end{fullwidth}
}

% }}}

\newcommand{\hlred}[1]{\textcolor{Maroon}{#1}} % Print text in maroon

\newcommand{\hangp}[1]{\makebox[0pt][r]{(}#1\makebox[0pt][l]{)}} % New command to create parentheses around text in tables which take up no horizontal space - this improves column spacing
\newcommand{\hangstar}{\makebox[0pt][l]{*}} % New command to create asterisks in tables which take up no horizontal space - this improves column spacing

\newcommand{\monthyear}{\ifcase\month\or January\or February\or March\or 
April\or May\or June\or July\or August\or September\or October\or 
November\or December\fi\space\number\year} % A command to print the current month and year

\newcommand{\openepigraph}[2]{ % This block sets up a command for printing an epigraph with 2 arguments - the quote and the author
\begin{fullwidth}
\sffamily\large
\begin{doublespace}
\noindent\allcaps{#1}\\ % The quote
\noindent\allcaps{#2} % The author
\end{doublespace}
\end{fullwidth}
}

\newcommand{\blankpage}{\newpage\hbox{}\thispagestyle{empty}\newpage} % Command to insert a blank page

\newcommand{\hangleft}[1]{\makebox[0pt][r]{#1}} % Used for printing commands in the index, moves the slash left so the command name aligns with the rest of the text in the index 
\newcommand{\hairsp}{\hspace{1pt}} % Command to print a very short space
\newcommand{\ie}{\textit{i.\hairsp{}e.}\xspace} % Command to print i.e.
\newcommand{\eg}{\textit{e.\hairsp{}g.}\xspace} % Command to print e.g.
\newcommand{\na}{\quad--} % Used in tables for N/A cells
\newcommand{\measure}[3]{#1/#2$\times$\unit[#3]{pc}} % Typesets the font size, leading, and measure in the form of: 10/12x26 pc.
\newcommand{\tuftebs}{\symbol{'134}} % Command to print a backslash in tt type in OT1/T1

% }}}

% }}}

% BOOK META-INFORMATION {{{
%---------------------------------------------------------------------------

% https://tex.stackexchange.com/questions/179839/tufte-book-title-page-customization

\newsavebox{\logo}
\savebox{\logo}{\includegraphics[height=3cm]{logo_unito.png}}

% Title of the book
\title[TWAS: Bridging the Gap between Genome, Transcriptome and 
Disease]{Transcriptome-Wide Association Studies: Bridging the Gap 
between Genome, Transcriptome and Disease} 

\author[FM]
{
	Supervisor: Prof. Paolo Provero
	\hfill
	Candidate: Federico Marotta
}

\publisher
{
	\usebox{\logo} \par \vspace{1cm}
	\smallcaps{Università degli Studi di Torino} \par
	\smallcaps{Scuola di Medicina} \par
	\smallcaps{Dipartimento di Biotecnologie Molecolari e Scienze per la Salute} \par
	\smallcaps{Corso di Laurea in Biotecnologie}
} 

\date{Tesi di Laurea, 18 Luglio 2018}

% }}}

% FRONT MATTER {{{
%---------------------------------------------------------------------------

\begin{document}

\frontmatter

% EPIGRAPH
%---------------------------------------------------------------------------

%\thispagestyle{empty}
%\openepigraph{The public is more familiar with bad design than good 
%design. It is, in effect, conditioned to prefer bad design, because 
%that is what it lives with. The new becomes threatening, the old 
%reassuring.}{Paul Rand, {\itshape Design, Form, and Chaos}}
%\vfill
%\openepigraph{A designer knows that he has achieved perfection not when 
%there is nothing left to add, but when there is nothing left to take 
%away.}{Antoine de Saint-Exup\'{e}ry}
%\vfill
%\openepigraph{\ldots the designer of a new system must not only be the 
%implementor and the first large-scale user; the designer should also 
%write the first user manual\ldots If I had not participated fully in 
%all these activities, literally hundreds of improvements would never 
%have been made, because I would never have thought of them or perceived 
%why they were important.}{Donald E. Knuth}

% TITLE PAGE
%---------------------------------------------------------------------------

% 1. the default way
%\maketitle % Print the title page

% 2. The custom title page
\makeatletter
\renewcommand{\maketitlepage}{%
\begingroup%
\begin{fullwidth}
\begin{centering}
\setlength{\parindent}{0pt}

\fontsize{13}{13}\selectfont\textit{\@publisher}\par

\null\vspace{\stretch{1}}

{\fontsize{22}{40}\selectfont\@title\par}

\vspace{\stretch{1}}

{\fontsize{16}{16}\selectfont\textit{\@author}\par}

\vspace{\stretch{1}}\null

{\fontsize{13}{13}\selectfont\textsf{\smallcaps{\@date}}\par}

\thispagestyle{empty}
\end{centering}
\end{fullwidth}
\endgroup
}
\makeatother
\maketitle % Print the title page

% COPYRIGHT PAGE
%---------------------------------------------------------------------------

\newpage
%\begin{fullwidth}
~\vfill
\thispagestyle{empty}
\setlength{\parindent}{0pt}
\setlength{\parskip}{\baselineskip}

\ccbysa

This work is licensed under a 
\href{http://creativecommons.org/licenses/by-sa/4.0/} {Creative Commons 
	Attribution-ShareAlike 4.0 International License}, and typeset with 
\href{ttps://www.latex-project.org/}{\LaTeX} using the 
\href{https://tufte-latex.github.io/tufte-latex/}{Tufte-LaTeX} class.

%\par\smallcaps{Published by \thanklesspublisher}
%\par\smallcaps{tufte-latex.googlecode.com}
%\par Licensed under the Apache License, Version 2.0 (the ``License''); 
%you may not use this file except in compliance with the License. You 
%may obtain a copy of the License at 
%\url{http://www.apache.org/licenses/LICENSE-2.0}. Unless required by 
%applicable law or agreed to in writing, software distributed under the 
%License is distributed on an \smallcaps{``AS IS'' BASIS, WITHOUT 
%WARRANTIES OR CONDITIONS OF ANY KIND}, either express or implied. See 
%the License for the specific language governing permissions and 
%limitations under the License.\index{license}
%\par\textit{First printing, \monthyear}
%\end{fullwidth}

% TABLES OF CONTENTS, FIGURES, TABLES
%---------------------------------------------------------------------------

% Decrease line spacing so that the toc fits in one page.
% https://tex.stackexchange.com/questions/56546/how-to-change-spaces-between-items-in-table-of-contents
\renewcommand{\baselinestretch}{0.75}\normalsize
\tableofcontents % Print the table of contents
\listoffigures % Print a list of figures
\addcontentsline{toc}{chapter}{List of Figures}
\listoftables % Print a list of tables
\addcontentsline{toc}{chapter}{List of Tables}
\listoftodos
\renewcommand{\baselinestretch}{1.0}\normalsize

% DEDICATION PAGE
%---------------------------------------------------------------------------

\cleardoublepage
\null\vspace{\stretch{1}}
\begin{doublespace}
\noindent\fontsize{15}{15}\selectfont\itshape
\nohyphenation

\enquote{The harmony of the world is made manifest in Form and Number, 
	and the heart and soul and all the poetry of Natural Philosophy are 
	embodied in the concept of mathematical beauty.}\\
---D'Arcy Wentworth Thompson, \textnormal{On Growth and Form}

\vspace{\stretch{1}}

\noindent\fontsize{15}{15}\selectfont\itshape

I would like to thank my supervisor, Prof. Paolo Provero, and all the 
former and current members of the Computational Biology Unit whom I have 
met: Elisa Mariella, Davide Marnetto, Elena Grassi, Ugo Ala, Alessandro 
Lussana, Stefano Gilotto. Their help has often been invaluable and they 
have never failed to provide a stimulating environment for me to work 
in.

\end{doublespace}
\vspace{\stretch{1}}\null

% ABSTRACT
%---------------------------------------------------------------------------

\cleardoublepage
\chapter{Abstract} % "\chapter*" would leave out this chapter from the table of contents

Understanding how genetic variation among individuals can influence the 
manifestation of complex diseases, which stem from the interaction of 
many genes with each other and with the environment, is a relevant 
problem in medicine. Since the first genome-wide association study (or 
GWAS) was conducted in 2005, tens of thousands of SNP-trait associations 
have been reported, shedding light on the at least partly genetic roots 
of many diseases; most of such associations, however, do not provide 
much predictive value and are difficult to explain. Expression 
quantitative trait loci (eQTL) mapping, which identifies loci that 
influence gene expression, is a possible step towards a better 
understanding of the relationship between genetic variation and 
phenotypic trait, using gene expression as a proxy for the trait. 
Recently, a new approach has been devised which goes one step further 
and aims to directly find associations between the expression of each 
gene and a given trait by combining GWAS and eQTL data. In one of their 
versions, these \enquote{transcriptome-wide association studies} (or 
TWAS) are performed in two phases, the first being the prediction of the 
genetic component of gene expression of the individuals in a GWAS cohort 
using reference transcriptome data, and the second being the evaluation 
of the association between predicted expression and trait in those 
individuals. On the whole, TWAS are powerful statistical methods to find 
associations between gene expression and complex phenotypic traits; 
while they can help in making sense of GWAS results, they also can find 
novel associations, pointing at potential candidate genes for further 
analysis: as such, their contribute to the characterisation of the 
relationship between genome and phenotype is substantial. After an 
introduction, the focus of the first part of this thesis will be a 
method to leverage individual-level data in order to detect genes 
associated with disease traits. The second part shall deal with how, 
conveniently, a TWAS can be performed starting only from the summary 
association statistics and the summary LD information of a GWAS. In the 
third part, we will discuss the advantages of integrating epigenetic 
markers in a TWA study and see an application to schizophrenia. 

% }}}

% MAIN MATTER {{{
%---------------------------------------------------------------------------

\mainmatter

\subfile{chapters/introduction.tex}

\subfile{chapters/methods.tex}

\subfile{chapters/gamazon2015.tex}

\subfile{chapters/gusev2016.tex}

\subfile{chapters/gusev2018.tex}

% }}}

% BACK MATTER {{{
%---------------------------------------------------------------------------

\backmatter

% BIBLIOGRAPHY
%---------------------------------------------------------------------------

\printbibheading[heading=bibintoc,title=References]
\printbibliography[heading=subbibliography,keyword=TWAS,title=Main 
Articles]
\printbibliography[heading=subbibliography,notkeyword=TWAS,title=Further 
Reading]
\nocite{*}

% APPENDIX
%---------------------------------------------------------------------------

%\appendix

% INDEX
%---------------------------------------------------------------------------

%\printindex % Print the index at the very end of the document (remember 
%to uncomment \makeindex if you use the index)

\end{document}

% }}}
